%!TEX root = ../main.tex

\section{Introduction}

Large language models (LLMs) are extremely performant on a wide range of natural language tasks, but they require enormous amounts of compute to train~\citep{openai2023gpt4tr, anthropicclaudeblog}.
As such, there is growing interest in building strong moderate-sized models, such as LLaMA~\citep{touvron2023llama,touvron2023llama2}, MPT~\citep{MosaicML2023Introducing}, and Falcon~\citep{falcon40b}, that allow for efficient inference and fine-tuning.
These LLMs are available in varied sizes suited for different use cases, but training each individual model from scratch---even the smallest billion-parameter models---requires substantial computational resources that are cost-prohibitive for most organizations. 
In this work, we seek to address the following question: 
\begin{center}
\vspace{-5pt}
    \emph{Can we produce a smaller, general-purpose, and competitive LLM by leveraging existing pre-trained LLMs, while using much less compute than training one from scratch?}
\vspace{-5pt}
\end{center}

We explore structured pruning as a means to achieve this goal. Pruning is commonly viewed as a solution for  compressing task-specific models~\citep{han2016deep,li2016pruning,lagunas2021block, xia-etal-2022-structured,kurtic2023ziplm}, removing redundant parameters and accelerating inference without sacrificing task performance. 
However, 
for general-purpose LLMs,
pruning inevitably results in performance degradation compared to original models~\citep{frantar2023sparsegpt,sun2023simple,ma2023llm}, especially when without significant compute invested post-pruning. In this work, we use pruning as an effective approach for developing smaller yet competitive LLMs that require only a fraction of the training compute compared to training them from scratch. 


We identify two key technical challenges in this problem. First, how can we decide on final pruned architectures that are strong in performance and efficient for inference? 
Existing structured pruning techniques for LLMs~\citep{xia-etal-2022-structured, ma2023llm} do not specify targeted structures and lead to suboptimal pruned models in terms of performance and inference speed (\Cref{tab:uniform} and \Cref{app:compare_to_ma}). 
Second, how can we continue pre-training the pruned model to reach desired performance? 
We observe that training using
the original pre-training data %
leads to imbalanced rates of loss reduction across different domains,
compared to when training such models from scratch.
This indicates that the pruned model retains varying levels of knowledge for different domains (e.g., GitHub vs. C4) and 
simply using the pre-training domain proportion results in an inefficient use of  data 
(\Cref{fig:dynamiclossdiff}). 
To address these issues, we propose ``{\method{}}'', an algorithm consisting of the following two components:
\begin{wrapfigure}[16]{r}{0.5\textwidth}
        \vspace{-0.9em}
        \centering
        \includegraphics[width=0.5\textwidth]{figures/teaserwlegend.pdf}
        \caption{
        \ours{}-2.7B surpasses a series of open-source models at a similar scale and only requires 1/32 (3\%) of budget to achieve on-par performance with OpenLLaMA-3B-v2. % 
    %    \protect\footnotemark
        }
        \vspace{-4em}
        \label{fig:teaser}
    \end{wrapfigure}
    % \footnotetext{\rebuttal{Note that the source model of Sheared-LLaMA is LLaMA2-7B, which is pretrained with 2T tokens.}}
\vspace{-1.5em}
\begin{itemize}[leftmargin=*,labelindent=1em]
    \item We propose a novel pruning algorithm, dubbed \textit{targeted structured pruning}, which prunes a source model to a specified target architecture. 
    The target architecture is determined by leveraging the configurations of existing pre-trained models. 
    Our pruning approach searches for substructures within the source model that maximally preserve performance while adhering to the given constraints. 
    \item 
    We devise a \textit{dynamic batch loading} algorithm that 
    loads training data from each domain in proportion to its rate of loss reduction, thereby 
    making an efficient use of the data and  
    accelerating the overall performance improvement. 
\end{itemize}

We demonstrate the efficacy of our proposed method by pruning a LLaMA2-7B model~\citep{touvron2023llama2} into two smaller LLMs:
\ours{}-1.3B and 
\ours{}-2.7B. 
Despite using only 50 billion addtional tokens (i.e., 5\% of OpenLLaMA's pre-training budget) for pruning and continued pre-training,
\ours{}-1.3B and \ours{}-2.7B 
outperform other popular LLMs at similar scales, including Pythia~\citep{biderman2023pythia}, INCITE~\citep{together2023redpajama}, and OpenLLaMA~\citep{openlm2023openllama}, on 11 representative downstream tasks (\Cref{fig:teaser}; commonsense, reading comprehension, and world knowledge) and instruction tuning for open-ended generation.
Additionally, the downstream performance trajectory suggests that further training the pruned model with more tokens would result in even greater gains.
While we only conduct experiments with up to 7B parameter models, our \method{} algorithm is highly generalizable and can be extended to large language models of any size in future work.
